\reviewercomment{The reviewer asks whether nutrient enrichment is being treated too simply as a direct proxy for stronger pairwise interaction strength. They note that enrichment can also alter carrying capacities, metabolic rates, pH feedbacks, and the balance of competition and facilitation, and suggest reframing the result in terms of interaction intensity, environmental feedbacks, and network restructuring rather than a monotonic increase in effective pairwise coefficients.}

We agree that nutrient enrichment should not be interpreted as a direct, monotonic proxy for stronger effective pairwise Lotka--Volterra coefficients. In the revised framing, the gLV parameter $\mu$ is a phenomenological control on average competitive suppression in the model, whereas the nutrient manipulation in the experiments is treated as a perturbation that changes net interaction intensity and environmental feedbacks, measured operationally through invasion resistance and coalescence outcomes. This distinction is important because the experimental nutrient gradient can affect multiple ecological processes at once, including resource competition, growth yield, environmental modification, and changes in community composition.

The data support a more nuanced statement. Pairwise failed-invasion frequency increased across the nutrient gradient from $2 \pm 1\%$ in Nutr$-$ to $33 \pm 4\%$ in Base and $48 \pm 4\%$ in Nutr$+$ (Fig.~4B; Supplementary Figs.~7--9), showing that invasion resistance increased under nutrient enrichment. Coalescence outcomes shifted in the same direction, with Dominance increasing from 39\% in Nutr$-$ to 65\% in Base and 76\% in Nutr$+$, while Mixture decreased from 53\% to 4\% and 6\% (Fig.~4D). However, we do not interpret these trends as evidence that external nutrient supply simply rescales all pairwise coefficients. Consistent with the reviewer's point, the manuscript now states that the mapping between nutrient concentration and $\mu$ is approximate and that the empirical signal may reflect denser resource competition, increased metabolic activity, pH-mediated environmental modification, and shifts in competition--facilitation balance.

\begin{figure}[H]
\centering
\includegraphics[width=0.95\textwidth]{figures/r2_q1_nutrient_interaction_feedback.pdf}
\caption{\textbf{Response Fig.~R2-Q1. Nutrient enrichment changes invasion resistance, coalescence outcomes, and environmental-feedback structure.} \textbf{A,} Failed invasion frequency and Dominance both increase across the nutrient gradient, but the failed-invasion assay is an operational measure of net invasion resistance rather than a direct estimate of a single biochemical mechanism. \textbf{B,} Nutrient enrichment is accompanied by increased dominant-ASV abundance and decreased parental richness, indicating network restructuring. \textbf{C,} In acid--alk coalescence events, the acidic parent wins strongly in Nutr$+$ but not significantly in Base, supporting pH modification as one possible contributor to winner identity in the nutrient-rich regime.}
\label{fig:response_r2_q1}
\end{figure}

The pH results are a concrete example of this broader interpretation. Dominant taxa include strong acidifiers and alkalizers, and dominant pH-modifying species predict parental community pH (Supplementary Figs.~17--18). In the strict acidic--alkaline subset used in Extended Data Fig.~7, where one parental community had pH $<6.5$ and the other had pH $>7.5$, the acidic parental community dominated 29/32 Nutr$+$ events (90.6\%), whereas the Base signal was weaker (23/41 events, 56.1\%). Thus, in Nutr$+$, environmental modification by dominant taxa may contribute to winner identity, while Base is better described as a more collective regime in which dominant-species predictability is weaker ($R^2 = 0.11$ in Base versus $R^2 = 0.49$ in Nutr$+$; Fig.~5C). This strict high-contrast analysis supports pH-mediated winner-direction bias, but it does not by itself explain the full Dominance pattern across all coalescence events.

We have therefore narrowed the claim from ``nutrient enrichment directly increases interaction strength'' to ``nutrient enrichment increases operational interaction intensity and environmental feedbacks in our assay.'' The revised manuscript distinguishes the model parameter $\mu$ from the experimental nutrient gradient, defines community-level selection as origin-correlated persistence rather than a single biochemical mechanism, and explicitly acknowledges that environmentally mediated interactions and network restructuring can contribute to the observed transition. We also propose adding the reviewer-suggested context on environmentally mediated assembly and consumer-resource mappings, citing Goldford et al., Estrela et al., and Duan et al. \citep{Goldford2018,Estrela2021,DuanPawar2025}.

\mschange{Manuscript change: Results and Supplementary Methods now frame the gLV model as phenomenological and state that $\alpha_{ij}$ can absorb net inhibitory effects from direct competition and indirect mechanisms such as pH modification, without identifying the biochemical mechanism. The Discussion now states that the nutrient-to-$\mu$ mapping is approximate and that enrichment can alter denser resource competition, metabolic activity, pH shifts, and competition--facilitation balance. We additionally propose final wording changes in the Results subsection title, Fig.~4 caption, Introduction, and Discussion to replace direct ``nutrient-dependent interaction strength'' language with ``nutrient-dependent interaction intensity and environmental feedbacks,'' as detailed in \texttt{proposed\_v5\_patches.md}.}
